\documentclass[11pt,a4paper]{article}
\usepackage[utf8]{inputenc}
\usepackage[spanish]{babel}
\usepackage{amsmath, amssymb}
\usepackage{geometry}
\geometry{top=2.5cm, bottom=2.5cm, left=2.5cm, right=2.5cm}

\begin{document}

\section*{Piense, dialogue y comparta}

\begin{itemize}
    \item[\textbf{1.}] Usted está en un carnaval jugando el juego ``Pégale a la campana'', como se muestra en la figura TP2.1. El objetivo es pegar en el extremo de la palanca con un martillo, enviando un objeto duro hacia arriba a lo largo de la pista vertical sin fricción con el fin de que le pegue a una campana en la parte superior. Mostrándose fuera de control para la multitud, usted pegó en la palanca varias veces en una fila de tal manera que el objeto se eleva una altura $h = 4.50 \text{ m}$ y apenas toca la campana, que hace un sonido de zumbido suave. Ahora, para impresionar realmente a la multitud, mueve el martillo con un movimiento poderoso, pega en la palanca y proyecta el objeto hacia arriba con el doble de la velocidad inicial de sus demostraciones anteriores. Pero usted no sabía que, en las demostraciones anteriores, la campana se soltó y se deslizó a un lado, por lo que en esta demostración el objeto se desvía de la campana y se proyecta hacia arriba en el aire. ¿Cuál es el intervalo de tiempo total entre que el objeto comienza su movimiento ascendente y después regresa a la tierra a un lado del aparato?
    \vspace*{9cm}

    \item[\textbf{2.}] Su grupo está en la cima de un acantilado de altura $h = 75.0 \text{ m}$. En el fondo del acantilado hay un ojo de agua. Usted divide al grupo en dos. Un miembro de la primera mitad de los voluntarios del grupo deja caer una piedra a partir del reposo esta cae hacia abajo en el agua. Un miembro de la segunda mitad de los voluntarios del grupo, después de que ha pasado algún intervalo de tiempo desde que se cayó la primera piedra, lanza una segunda piedra hacia abajo para que ambas lleguen al agua al mismo tiempo. Prueba la ejecución escuchando una sola salpicada hecha por las piedras al mismo tiempo que pegan en el agua. (a) Si la segunda piedra se lanza $1.00 \text{ s}$ después de que se lanza la primera, ¿con qué rapidez se debe lanzar la segunda piedra? (b) Si o más rápido que cualquiera en su grupo puede lanzar la piedra es $40.0 \text{ m/s}$, ¿cuál es el intervalo de tiempo más largo que puede pasar entre la liberación de las piedras para que se escuche una sola salpicada? (c) Si no hay límite en cuanto a la velocidad con que se puede lanzar la roca, ¿cuál es el intervalo de tiempo más largo que puede pasar entre la liberación de las piedras para que se escuche una sola salpicada?
    \vspace*{12cm}

    \item[\textbf{3.}] \textbf{ACTIVIDAD} Haga que su compañero sostenga una regla verticalmente con el extremo cero en la parte inferior. Coloque el dedo y el pulgar abiertos en la posición cero. Sin previo aviso, su compañero debe soltar la regla y usted debe agarrarla tan pronto como usted la ve en movimiento. Desde la posición de su dedo en la regla, determine su tiempo de reacción. Repita el experimento varias veces para estimar la incertidumbre en su tiempo de reacción. Haga que cada miembro de su grupo atrape a la regla y compare sus tiempos de reacción.
    \vspace*{6cm}

    \item[\textbf{4.}] \textbf{ACTIVIDAD} El Acela es un tren eléctrico en la ruta Washington-Nueva York-Boston y transporta pasajeros a $170 \text{ mi/h}$. En la figura TP2.4 se muestra una gráfica velocidad-tiempo para el Acela. (a) Describa el movimiento del tren en cada intervalo de tiempo sucesivo. (b) Encuentre la aceleración máxima positiva del tren en la gráfica del movimiento. (c) Encuentre el desplazamiento del tren, en millas, entre $t = 0$ y $t = 200 \text{ s}$.
    \vspace*{9cm}
\end{itemize}

\newpage

\section*{Problemas}

\subsection*{SECCIÓN 2.1 Posición, velocidad y rapidez de una partícula}

\begin{itemize}
    \item[\textbf{1.}] La rapidez de un impulso nervioso en el cuerpo humano es de aproximadamente $100 \text{ m/s}$. Si su dedo del pie tropieza accidentalmente en la oscuridad, estime el tiempo que tarda el impulso nervioso en viajar a su cerebro.
    \vspace*{5cm}

    \item[\textbf{2.}] Una partícula se mueve de acuerdo con la ecuación $x = 10t^2$, donde $x$ está en metros y $t$ en segundos. (a) Encuentre la velocidad promedio para el intervalo de tiempo de $2.00 \text{ s}$ a $3.00 \text{ s}$. (b) Encuentre la velocidad promedio para el intervalo de tiempo de $2.00 \text{ s}$ a $2.10 \text{ s}$.
    \vspace*{6cm}

    \item[\textbf{3.}] La posición de un automóvil se observó en varios momentos; los resultados se resumen en la tabla siguiente. Encuentre la velocidad promedio del auto para (a) el primer segundo, (b) los últimos $3 \text{ s}$ y c) todo el periodo de observación.
    \begin{center}
    \begin{tabular}{lcccccc}
    \hline
    $t$ (s) & 0 & 1.0 & 2.0 & 3.0 & 4.0 & 5.0 \\
    $x$ (m) & 0 & 2.3 & 9.2 & 20.7 & 36.8 & 57.5 \\
    \hline
    \end{tabular}
    \end{center}
    \vspace*{6cm}
\end{itemize}

\newpage

\subsection*{SECCIÓN 2.2 Velocidad y rapidez instantáneas}

\begin{itemize}
    \item[\textbf{4.}] Una atleta parte desde un extremo de una piscina de longitud $L$ en $t = 0$ y llega en el otro extremo en el tiempo $t_1$. Nada de regreso y llega a la posición de partida en el tiempo $t_2$. Si ella está nadando inicialmente en la dirección $x$ positiva, determine sus velocidades promedio simbólicamente en (a) la primera parte del nado, (b) la segunda mitad del nado y (c) el recorrido redondo. (d) ¿Cuál es su rapidez promedio para el recorrido redondo?
    \vspace*{7cm}

    \item[\textbf{5.}] En la figura P2.5 se muestra una gráfica posición-tiempo para una partícula que se mueve a lo largo del eje $x$. (a) Encuentre la velocidad promedio en el intervalo de tiempo $t = 1.50 \text{ s}$ a $t = 4.00 \text{ s}$. (b) Determine la velocidad instantánea en $t = 2.00 \text{ s}$ al medir la pendiente de la recta tangente que se muestra en la gráfica. (c) ¿En qué valor de $t$ la velocidad es cero?
    \vspace*{7cm}
\end{itemize}

\newpage

\subsection*{SECCIÓN 2.3 Modelo de análisis: Partícula bajo velocidad constante}

\begin{itemize}
    \item[\textbf{6.}] Un auto viaja a lo largo de una línea recta a una rapidez constante de $60.0 \text{ mi/h}$ una distancia $d$ y luego otra distancia $d$ en la misma dirección con otra rapidez constante. El promedio de velocidad para todo el viaje es de $30.0 \text{ mi/h}$. (a) ¿Cuál es la rapidez constante con la que el auto se trasladó durante la segunda distancia $d$? (b) \textbf{¿Qué pasaría si?} Supongamos que la segunda distancia $d$ se viajó en dirección opuesta; olvidó algo y tuvo que volver a casa con la misma rapidez constante que encontró en el inciso (a). ¿Cuál es la velocidad promedio para este viaje? (c) ¿Cuál es la rapidez promedio para este nuevo viaje?
    \vspace*{7cm}

    \item[\textbf{7.}] Una persona que hace un viaje, conduce con una rapidez constante de $89.5 \text{ km/h}$, excepto por una parada de descanso de $22.0 \text{ min}$. Si la rapidez promedio de la persona es de $77.8 \text{ km/h}$, (a) ¿cuánto tiempo invierte la persona en el viaje y (b) qué tan lejos llegará?
    \vspace*{6cm}
\end{itemize}

\newpage

\subsection*{SECCIÓN 2.5 Aceleración}

\begin{itemize}
    \item[\textbf{8.}] Una niña rueda una canica sobre una pista con dobleces que mide $100 \text{ cm}$ de largo, como se muestra en la figura P2.8. Use $x$ para representar la posición de la canica a lo largo de la pista. En las secciones horizontales de $x = 0$ a $x = 20 \text{ cm}$ y de $x = 40 \text{ cm}$ a $x = 60 \text{ cm}$, la canica rueda con rapidez constante. En las secciones de pendiente, la rapidez de la canica cambia de manera uniforme. En los lugares donde la pendiente cambia, la canica permanece en la pista y no experimenta cambios súbitos en rapidez. La niña da a la canica cierta rapidez inicial en $x = 0$ y $t = 0$ y luego la observa rodar a $x = 90 \text{ cm}$ donde regresa, y eventualmente regresa a $x = 0$ con la misma rapidez con la que al inicio la niña la liberó. Trace gráficas de $x$ en función de $t$, $v_x$ en función de $t$ y $a_x$ en función de $t$, alineadas verticalmente con sus ejes de tiempo idénticos, para mostrar el movimiento de la canica. No podrá colocar números distintos a cero en el eje horizontal o en los ejes de velocidad o aceleración, solo muestre las formas correctas en las gráficas.
    \vspace*{10cm}

    \item[\textbf{9.}] La figura P2.9 muestra una gráfica de $v_x$ en función de $t$ para el movimiento de un motociclista mientras parte del reposo y se mueve a lo largo del camino en línea recta. (a) Encuentre la aceleración promedio para el intervalo de tiempo $t = 0$ a $t = 6.00 \text{ s}$. (b) Estime el tiempo en que la aceleración tiene su mayor valor positivo y el valor de la aceleración en dicho instante. (c) ¿Cuándo la aceleración es cero? (d) Estime el máximo valor negativo de la aceleración y el tiempo en el que ocurre.
    \vspace*{7cm}

    \item[\textbf{10.}] (a) Use los datos del problema 3 para construir una gráfica continua de posición en función del tiempo. (b) Con la construcción de rectas tangentes a la curva $x(t)$, encuentre la velocidad instantánea del automóvil en varios instantes. (c) Grafique la velocidad instantánea en función del tiempo y, a partir de la gráfica, determine la aceleración promedio del automóvil. (d) ¿Cuál fue la velocidad inicial del automóvil?
    \vspace*{10cm}

    \item[\textbf{11.}] Una partícula parte del reposo y acelera como se muestra en la figura P2.11. Determine (a) la rapidez de la partícula en $t = 10.0 \text{ s}$ y en $t = 20.0 \text{ s}$ y (b) la distancia recorrida en los primeros $20.0 \text{ s}$.
    \vspace*{7cm}
\end{itemize}

\newpage

\subsection*{SECCIÓN 2.6 Diagramas de movimiento}

\begin{itemize}
    \item[\textbf{12.}] Dibuje diagramas de movimiento para (a) un objeto que se mueve a la derecha con rapidez constante, (b) un objeto que se mueve a la derecha y está aumentando su rapidez de modo constante, (c) un objeto que se mueve a la derecha y está disminuyendo su rapidez de modo constante, (d) un objeto que se mueve a la izquierda y aumenta su rapidez en proporción constante y (e) un objeto que se mueve a la izquierda y frena a rapidez constante. (f) ¿Cómo modificaría su dibujo si los cambios en rapidez no fuesen uniformes; esto es, si la rapidez no cambiara de modo constante?
    \vspace*{8cm}

    \item[\textbf{13.}] Cada una de las fotografías estroboscópicas (a), (b) y (c) en la figura P2.13 se tomó de un solo disco que se mueve hacia la derecha, que se considera como la dirección positiva. Dentro de cada fotografía el intervalo de tiempo entre imágenes es constante. Para cada fotografía prepare gráficas de $x$ en función de $t$, $v_x$ en función de $t$ y $a_x$ en función de $t$, alineadas verticalmente con sus ejes de tiempo idénticos, para mostrar el movimiento del disco. No podrá colocar números distintos de cero sobre los ejes, sino mostrar los tamaños relativos correctos sobre las gráficas.
    \vspace*{10cm}
\end{itemize}

\newpage

\subsection*{SECCIÓN 2.7 Modelo de análisis: Partícula bajo aceleración constante}

\begin{itemize}
    \item[\textbf{14.}] Un electrón en un tubo de rayos catódicos acelera uniformemente desde una rapidez de $2.00 \times 10^4 \text{ m/s}$ a $6.00 \times 10^6 \text{ m/s}$ en $1.50 \text{ cm}$. (a) ¿En qué intervalo de tiempo el electrón recorre estos $1.50 \text{ cm}$? (b) ¿Cuál es su aceleración?
    \vspace*{6cm}

    \item[\textbf{15.}] Considere una porción de aire en un tubo recto que se mueve con una aceleración constante de $-4.00 \text{ m/s}^2$ y tiene una velocidad de $13.0 \text{ m/s}$ a las 10:05:00 a.m. (a) ¿Cuál es su velocidad a las 10:05:01 a.m.? (b) ¿A las 10:05:04 a.m.? (c) ¿A las 10:04:59 a.m.? (d) Describa la forma de una gráfica de velocidad en función del tiempo para esta porción de aire. (e) Argumente a favor o en contra del enunciado: ``Conocer un solo valor de la aceleración constante de un objeto es como conocer toda una lista de valores para su velocidad''.
    \vspace*{7cm}

    \item[\textbf{16.}] En el ejemplo 2.7, se investigó un aterrizaje del jet en un portaaviones. En una maniobra posterior, el jet viene para un aterrizaje en tierra firme con una rapidez de $100 \text{ m/s}$, y su aceleración puede tener una magnitud máxima de $5.00 \text{ m/s}^2$, hasta alcanzar el reposo. (a) Desde el momento en que el avión toca la pista de aterrizaje, ¿cuál es el intervalo de tiempo mínimo necesario antes de que pueda alcanzar el reposo? (b) ¿Puede que el jet aterrice en un pequeño aeropuerto de la isla tropical, donde la pista es de $0.800 \text{ km}$ de largo? (c) Explique su respuesta.
    \vspace*{6cm}

    \item[\textbf{17.}] Un objeto que se mueve con aceleración uniforme tiene una velocidad de $12.0 \text{ cm/s}$ en la dirección $x$ positiva cuando su coordenada $x$ es $3.00 \text{ cm}$. Si su coordenada $x$, $2.00 \text{ s}$ después es $-5.00 \text{ cm}$, ¿cuál es su aceleración?
    \vspace*{6cm}
\end{itemize}

\newpage

\begin{itemize}
    \item[\textbf{18.}] Resuelva el ejemplo 2.8 mediante un método gráfico. En la misma gráfica trace posición en función del tiempo para el automóvil y el oficial de policía. De la intersección de las dos curvas lea el tiempo cuando el patrullero alcanza al automóvil.
    \vspace*{7cm}

    \item[\textbf{19.}] Un deslizador de longitud $\lambda$ en una pista de aire se mueve a través de una fotocelda fija. Una fotocelda (figura P2.19) es un dispositivo que mide el intervalo de tiempo durante el $\Delta t_d$ que bloquea al deslizador de un haz de luz infrarroja que pasa a través de la fotopuerta. La proporción $v_d = \lambda / \Delta t_d$ es la velocidad promedio del deslizador durante esta parte de su movimiento. Suponga que el deslizador se mueve con aceleración constante. (a) Argumente a favor o en contra de la idea de que $v_d$ es igual a la velocidad instantánea del deslizador cuando está a la mitad de la fotocelda en el espacio. (b) Argumente a favor o en contra de la idea de que $v_d$ es igual a la velocidad instantánea del deslizador cuando está a la mitad de la fotocelda en el tiempo.
    \vspace*{6cm}

    \item[\textbf{20.}] ¿Por qué no es posible la siguiente situación? A partir del reposo, un rinoceronte se mueve a $50.0 \text{ m}$ en una línea recta en $10.0 \text{ s}$. Su aceleración es constante durante todo el movimiento y su rapidez final es $8.00 \text{ m/s}$.
    \vspace*{5cm}

    \item[\textbf{21.}] Un deslizador de $12.4 \text{ cm}$ de longitud se mueve sobre una pista de aire con aceleración constante. Transcurre un intervalo de tiempo de $0.628 \text{ s}$ entre el momento en que su extremo frontal pasa un punto fijo \textcircled{A} a lo largo de la pista y el momento cuando su extremo trasero pasa este punto. A continuación, transcurre un intervalo de tiempo de $1.39 \text{ s}$ entre el momento en que el extremo trasero del deslizador pasa el punto \textcircled{A} y el momento en que el extremo frontal del deslizador pasa un segundo punto \textcircled{B} más lejos en la pista. Después de ello, transcurren $0.431 \text{ s}$ adicionales hasta que el extremo trasero del deslizador pasa el punto \textcircled{B}. (a) Encuentre la rapidez promedio del deslizador conforme pasa el punto \textcircled{A}. (b) Encuentre la aceleración del deslizador. (c) Explique cómo calcula la aceleración sin saber la distancia entre los puntos \textcircled{A} y \textcircled{B}.
    \vspace*{7cm}
\end{itemize}

\newpage

\begin{itemize}
    \item[\textbf{22.}] En el modelo de partícula bajo aceleración constante se identifican las variables y parámetros $v_{xi}$, $v_{xf}$, $a_x$, $t$ y $x_f - x_i$. De las ecuaciones en el modelo, ecuaciones 2.13-2.17, la primera no involucra $x_f - x_i$, la segunda y tercera no contienen $a_x$, la cuarta omite $v_{xf}$ y la última deja fuera $t$. De modo que, para completar el conjunto, debe haber una ecuación que no involucre $v_{xi}$. Dedúzcala a partir de las otras.
    \vspace*{6cm}

    \item[\textbf{23.}] En $t = 0$, un carro de juguete se pone a rodar en una pista recta con posición inicial de $15.00 \text{ cm}$, velocidad inicial de $-3.50 \text{ cm/s}$ y aceleración constante de $2.40 \text{ m/s}^2$. En el mismo momento, otro carro de juguete se pone a rodar en una pista adyacente con posición inicial de $10.0 \text{ cm}$, una velocidad inicial de $+5.50 \text{ cm/s}$ y aceleración constante cero. (a) ¿En qué tiempo, si hay alguno, los dos carros tienen iguales rapideces? (b) ¿Cuáles son sus rapideces en dicho tiempo? (c) ¿En qué tiempo(s), si hay alguno, los carros se rebasan mutuamente? (d) ¿Cuáles son sus ubicaciones en dicho tiempo? (e) Explique la diferencia entre la pregunta (a) y la pregunta (c) tan claramente como le sea posible.
    \vspace*{8cm}

    \item[\textbf{24.}] Usted está observando los postes a lo largo del lado de la carretera como se describe en el Imagine introductorio del capítulo. Usted ya se detuvo y midió la distancia entre los postes adyacentes de $40.0 \text{ m}$. Ahora está conduciendo de nuevo y ha activado el cronómetro del teléfono inteligente. Se inicia el cronómetro en $t = 0$ en cuanto pasa el poste 1. En el poste 2, el cronómetro indica $10.0 \text{ s}$. En el poste 3, el cronómetro indica $25.0 \text{ s}$. Su amigo le dice que estaba presionando el freno y ralentizando el auto uniformemente durante todo el intervalo de tiempo desde el poste 1 hasta el poste 3. (a) ¿cuál fue la aceleración del automóvil entre los postes 1 y 3? (b) ¿Cuál era la velocidad del auto en el poste 1? (c) Si el movimiento del auto continúa como se describe, ¿cuál es el número del último poste que pasa antes de que el auto alcance el reposo?
    \vspace*{8cm}
\end{itemize}

\newpage

\subsection*{SECCIÓN 2.8 Objetos en caída libre}

\begin{itemize}
    \item[\textbf{25.}] ¿Por qué no es posible la siguiente situación? Emily desafía a su amigo David a atrapar un billete de dólar del modo siguiente. Ella sostiene el billete verticalmente, como se muestra en la figura P2.25, con el centro del billete entre los dedos índice y pulgar de David, quien debe atrapar el billete después de que Emily lo libere sin mover su mano hacia abajo. Sin previo aviso, Emily libera el billete. David toma el billete sin mover la mano hacia abajo. El tiempo de reacción de David es igual al tiempo promedio de reacción humano.
    \vspace*{6cm}

    \item[\textbf{26.}] Un atacante en la base de la pared de un castillo de $3.65 \text{ m}$ de alto lanza una roca recta hacia arriba con una rapidez de $7.40 \text{ m/s}$ desde una altura de $1.55 \text{ m}$ sobre el suelo. (a) ¿La roca llegará a lo alto de la pared? (b) Si es así, ¿cuál es su rapidez en lo alto? Si no, ¿qué rapidez inicial debe tener para llegar a lo alto? (c) Encuentre el cambio en rapidez de una roca lanzada recta hacia abajo desde lo alto de la pared con una rapidez inicial de $7.40 \text{ m/s}$ y que se mueve entre los mismos dos puntos. (d) ¿El cambio en la rapidez de la roca que se mueve hacia abajo concuerda con la magnitud del cambio de rapidez de la roca que se mueve hacia arriba entre las mismas elevaciones? Explique físicamente por qué sí o por qué no concuerda.
    \vspace*{8cm}

    \item[\textbf{27.}] La altura de un helicóptero sobre el suelo está dada por $h = 3.00t^3$, donde $h$ está en metros y $t$ en segundos. Después de $2.00 \text{ s}$, el helicóptero libera una pequeña valija de correo. ¿Cuánto tiempo después de su liberación, la valija llega al suelo?
    \vspace*{6cm}

    \item[\textbf{28.}] Se lanza una pelota hacia arriba desde el suelo con una rapidez inicial de $25 \text{ m/s}$; en el mismo instante, se deja caer otra pelota desde un edificio de $15 \text{ m}$ de altura. ¿Después de cuánto tiempo estarán las pelotas a la misma altura por encima del suelo?
    \vspace*{6cm}
\end{itemize}

\newpage

\begin{itemize}
    \item[\textbf{29.}] Un estudiante lanza un conjunto de llaves verticalmente hacia arriba a su hermana de fraternidad, quien está en una ventana $4.00 \text{ m}$ arriba. La segunda estudiante atrapa las llaves $1.50 \text{ s}$ después. (a) ¿Con qué velocidad inicial se lanzaron las llaves? (b) ¿Cuál fue la velocidad de las llaves justo antes de ser atrapadas?
    \vspace*{6cm}

    \item[\textbf{30.}] Al tiempo $t = 0$, un estudiante lanza un juego de llaves verticalmente hacia arriba a su hermana de fraternidad, que está en una ventana a una distancia $h$ arriba. La segunda estudiante atrapa las llaves al tiempo $t$. (a) ¿Con qué velocidad inicial se lanzaron las llaves? (b) ¿Cuál era la velocidad de las llaves justo antes de que fueran capturadas?
    \vspace*{6cm}

    \item[\textbf{31.}] Usted ha sido contratado por la Fiscalía como un testigo experto en un caso de robo. El indiciado es acusado de robar un anillo de diamantes caro y de gran tamaño en su caja de una joyería. Un testigo del presunto crimen testificó que vio al acusado huir de la tienda, parar junto a un edificio de apartamentos, y aventar la caja hacia arriba a un cómplice que se asomaba de una ventana del cuarto piso. Cuando se capturó, el acusado no tenía la caja robada con él y reclamó inocencia. Cuando el testigo ratificó en el Tribunal sobre el lanzamiento de la caja acusó al cómplice, el abogado defensor argumentó que sería imposible aventar la caja hacia arriba tan alto para que llegue a la ventana en cuestión. La parte inferior de la ventana está $19.0 \text{ m}$ arriba de la acera. Usted ha montado una demostración en la que el juez le pide al acusado lanzar una pelota de béisbol horizontalmente tan rápido como pueda y se utiliza un dispositivo de radar para determinar que puede lanzar la pelota a $20 \text{ m/s}$. (a) ¿Qué testimonio puede proporcionar sobre la capacidad del acusado para lanzar la caja a la ventana en cuestión? (b) ¿Qué argumento podría hacer el abogado defensor sobre el proceso utilizado para desarrollar su testimonio de experto? ¿Cuál podría ser su argumento contrario? Desprecie cualquier efecto de la resistencia del aire en la caja.
    \vspace*{8cm}
\end{itemize}

\newpage

\subsection*{SECCIÓN 2.9 Ecuaciones cinemáticas deducidas del cálculo}

\begin{itemize}
    \item[\textbf{32.}] Un estudiante conduce un ciclomotor a lo largo de un camino recto como se describe en la gráfica velocidad en función del tiempo de la figura P2.32. Trace esta gráfica en una hoja de papel gráfico. (a) Directamente sobre su gráfica, trace una gráfica de la posición en función del tiempo y alinee las coordenadas de tiempo de las dos gráficas. (b) Trace una gráfica de la aceleración en función del tiempo directamente bajo de la gráfica velocidad-tiempo, y de nuevo alinee las coordenadas de tiempo. Sobre cada gráfica muestre los valores numéricos de $x$ y $a_x$ para todos los puntos de inflexión. (c) ¿Cuál es la aceleración en $t = 6.00 \text{ s}$? (d) Encuentre la posición (relativa al punto de partida) en $t = 6.00 \text{ s}$. (e) ¿Cuál es la posición final del ciclomotor en $t = 9.00 \text{ s}$?
    \vspace*{10cm}
\end{itemize}

\newpage

\subsection*{PROBLEMAS ADICIONALES}

\begin{itemize}
    \item[\textbf{33.}] Los ingenieros automotrices se refieren a la rapidez de cambio de la aceleración en el tiempo como el ``jalón''. Suponga que un objeto se mueve en una dimensión tal que su jalón $J$ es constante. (a) Determine expresiones para su aceleración $a_x(t)$, velocidad $v_{x}(t)$ y posición $x(t)$, dado que su aceleración, velocidad y posición iniciales son $a_{xi}$, $v_{xi}$ y $x_i$, respectivamente. (b) Muestre que $a_x^2 = a_{xi}^2 + 2J(v_x - v_{xi})$.
    \vspace*{7cm}

    \item[\textbf{34.}] En la figura 2.11b, el área bajo la curva velocidad en función del tiempo y entre el eje vertical y el tiempo $t$ (línea discontinua vertical) representa el desplazamiento. Como se muestra, esta área consiste en un rectángulo y un triángulo. (a) Calcule sus áreas y (b) explique cómo se compara la suma de las dos áreas con la expresión en el lado derecho de la ecuación 2.16.
    \vspace*{6cm}

    \item[\textbf{35.}] El saltamontes (\textit{Philaenus spumarius}) es supuestamente el mejor saltador en el reino animal. Para iniciar un salto, este insecto es capaz de acelerar a $4.00 \text{ km/s}^2$ en una distancia de $2.00 \text{ mm}$, ya que endereza sus especialmente adaptadas ``piernas saltadoras''. Suponga que la aceleración es constante. (a) Encuentre la velocidad ascendente con el que el insecto despega. (b) ¿En qué intervalo de tiempo se alcanza esta velocidad? (c) ¿A qué altura saltaría el insecto si la resistencia del aire fuera insignificante? La altura real a la que llega es aproximadamente de $70 \text{ cm}$, por lo que la resistencia del aire debe ser una fuerza importante en el salto del insecto.
    \vspace*{6cm}

    \item[\textbf{36.}] Una mujer informó haber caído $144$ metros del piso 17 de un edificio, quedando sobre una caja de metal del ventilador que aplastó a una profundidad de $18.0 \text{ pulg}$. Ella sufrió solo heridas leves. Ignorando la resistencia del aire, calcule (a) la rapidez de la mujer justo antes de que chocó con el ventilador y (b) su aceleración promedio mientras está en contacto con la caja. (c) Modele su aceleración como constante, calcule el intervalo de tiempo que se tardó en aplastar la caja.
    \vspace*{6cm}
\end{itemize}

\newpage

\begin{itemize}
    \item[\textbf{37.}] En $t = 0$, un atleta corriendo en una competencia de carrera en una pista recta larga, con una rapidez constante $v_1$, está atrás una distancia $d_1$ de un segundo atleta que corre con una rapidez constante $v_2$. (a) ¿En qué circunstancias el primer atleta es capaz de superar el segundo atleta? (b) Encuentre el tiempo $t$ en el que el primer atleta alcanza el segundo atleta, en términos de $d_1$, $v_1$ y $v_2$. (c) ¿A qué distancia mínima $d_2$ del líder atleta debe encontrarse la línea de meta para que el atleta que va detrás al menos pueda empatar en el primer lugar? Exprese $d_2$ en términos de $d_1$, $v_1$ y $v_2$ utilizando el resultado del inciso (b).
    \vspace*{6cm}

    \item[\textbf{38.}] ¿Por qué no es posible la siguiente situación? Un tren de carga viaja lentamente con una rapidez constante de $16.0 \text{ m/s}$. Detrás del tren de carga en la misma vía está un tren de pasajeros que viaja en la misma dirección a $40.0 \text{ m/s}$. Cuando la parte delantera del tren de pasajeros está a $58.5 \text{ m}$ de la parte trasera del tren de carga, el ingeniero del tren de pasajeros reconoce el peligro y frena su tren, haciendo que el tren se mueva con una aceleración $-3.00 \text{ m/s}^2$. Debido a la acción del ingeniero, los trenes no chocan.
    \vspace*{6cm}

    \item[\textbf{39.}] Hannah prueba su nuevo automóvil deportivo corriendo con Sam que es un experimentado corredor, pero Hannah deja la línea de partida $1.00 \text{ s}$ después que Sam. Sam se mueve con una aceleración constante de $3.50 \text{ m/s}^2$ y Hannah mantiene una aceleración de $4.90 \text{ m/s}^2$. Encuentre (a) el tiempo cuando Hannah alcanza a Sam, (b) la distancia que recorre antes de alcanzarlo y (c) las rapideces de ambos automóviles en el instante en que lo rebasa.
    \vspace*{6cm}

    \item[\textbf{40.}] Dos objetos, A y B, se conectan mediante una barra rígida que tiene longitud $L$. Los objetos se deslizan a lo largo de rieles guía perpendiculares como se muestra en la figura P2.40. Suponga que A se desliza hacia la izquierda con una rapidez constante $v$. (a) Encuentre la velocidad de B como una función del ángulo $\theta$. (b) Describa $v_B$ en relación con $v$. ¿Es $v_B$ siempre menor que $v$, mayor que $v$ o igual a $v$, o tiene alguna otra relación?
    \vspace*{6cm}

    \item[\textbf{41.}] Lisa llega de prisa a un andén del metro para encontrar que su tren está a punto de partir. Se detiene y mira los carros que pasan. Cada carro tiene una longitud de $8.60 \text{ m}$. El primero se mueve pasándola en $1.50 \text{ s}$ y el segundo en $1.10 \text{ s}$. Encuentre la aceleración constante del tren.
    \vspace*{6cm}
\end{itemize}

\newpage

\subsection*{PROBLEMAS DE DESAFÍO}

\begin{itemize}
    \item[\textbf{42.}] Dos varillas delgadas se sujetan en el interior de un anillo circular como se muestra en la figura P2.42. Una varilla de longitud $D$ es vertical, y la otra de longitud $L$ forma un ángulo $\theta$ con la horizontal. Las dos varillas y el anillo están en un plano vertical. Dos pequeñas perlas están libres para deslizarse sin fricción a lo largo de las varillas. (a) Si las dos perlas se liberan simultáneamente a partir del reposo desde las posiciones que se muestran, use su intuición y conjeture qué perla llega al fondo primero. (b) Halle una expresión para el intervalo de tiempo necesario para que la perla roja caiga del punto \textcircled{A} al punto \textcircled{C} en términos de $g$ y $D$. (c) Encuentre una expresión para el intervalo de tiempo requerido para la perla azul se deslice desde el punto \textcircled{B} al punto \textcircled{C} en términos de $g$, $L$ y $\theta$. (d) Demuestre que los dos intervalos de tiempo encontrados en los incisos (b) y (c) son iguales. \textit{Sugerencia:} ¿Cuál es el ángulo entre las cuerdas \textcircled{A}\textcircled{B} y \textcircled{B}\textcircled{C} del círculo? (e) ¿Estos resultados le sorprenden? Su conjetura intuitiva del inciso (a) ¿es correcta? Este problema se inspiró en un artículo de Thomas B. Greenslade, Jr., ``La paradoja de Galileo'', \textit{Phys. Teach.} \textbf{46}, 294 (mayo de 2008).
    \vspace*{8cm}

    \item[\textbf{43.}] En una carrera de $100 \text{ m}$ de mujeres, a Laura acelerando de manera uniforme, le toma $2.00 \text{ s}$ y Healan $3.00 \text{ s}$ alcanzar sus rapideces máximas, que mantienen cada una de ellas durante el resto de la carrera. Cruzan la línea de meta al mismo tiempo, estableciendo ambas un récord mundial de $10.4 \text{ s}$. (a) ¿Cuál es la aceleración de cada atleta? (b) ¿Cuáles son sus respectivas rapideces máximas? (c) ¿Qué atleta está por delante en la marca de $6.00 \text{ s}$, y por cuánto? (d) ¿Cuál es la distancia máxima en que Healan está detrás de Laura, y en qué tiempo eso ocurre?
    \vspace*{8cm}

    \item[\textbf{44.}] \textbf{Problema de repaso.} Usted está sentado en su auto en reposo en un semáforo con un ciclista en reposo junto a usted en el carril contiguo de bicicletas. Tan pronto como el semáforo se pone verde, un automóvil acelera a partir del reposo a $50.0 \text{ mi/h}$ con una aceleración constante de $9.00 \text{ mi/h/s}$ y después se mueve con una rapidez constante de $50.0 \text{ mi/h}$. Al mismo tiempo el ciclista acelera desde el reposo hasta $20.0 \text{ mi/h}$ con una aceleración constante de $13.0 \text{ mi/h/s}$ y después se mueve con rapidez constante de $20.00 \text{ mi/h}$. (a) ¿Durante qué intervalo de tiempo está la bicicleta por delante del automóvil? (b) ¿Para qué distancia máxima la bicicleta está delante del automóvil?
    \vspace*{8cm}
\end{itemize}

\end{document}
